\begin{table}[htbp]
\centering
\scriptsize
\caption{Mapa kanałów transmisji scenariusza BDP}
\label{tab:transmission-channels}
\begin{tabularx}{\textwidth}{@{}>{\raggedright\arraybackslash}p{0.16\textwidth}>{\raggedright\arraybackslash}p{0.32\textwidth}>{\raggedright\arraybackslash}p{0.18\textwidth}>{\raggedright\arraybackslash}X@{}}
\toprule
\textbf{Kanał} & \textbf{Mechanizm ekonomiczny} & \textbf{Wskaźniki} & \textbf{Interpretacja wyniku w wariancie centralnym} \\
\midrule
Płaca rezerwowa &
Transfer zwiększa dochód niezależny od pracy; część transferu podnosi płacę rezerwową przez parametr \(\lambda\). &
płaca rynkowa, płaca zatrudnionych, bezrobocie &
Kanał jest aktywny mechanicznie, ale płaca rynkowa reaguje silniej dopiero przy wysokim BDP; sam kanał nie tłumaczy lokalnego maksimum adopcji przy \texttt{--bdp 2000}. \\
\addlinespace
Popyt i dochód gospodarstw &
BDP zwiększa dochód gospodarstw domowych, lecz część impulsu może trafić do oszczędności, importu i bilansów zamiast do krajowego popytu finalnego. &
dochód HH, mediana konsumpcji, PKB, transfery społeczne &
Dochód gospodarstw rośnie, ale mediana konsumpcji i PKB reagują słabiej; kanał popytowy wspiera firmy tylko częściowo. \\
\addlinespace
Fiskus i dług &
Transfer zwiększa wydatki publiczne, deficyt oraz dług publiczny względem PKB. &
deficyt/PKB, dług/PKB, rentowność długu, obsługa długu &
To najsilniejszy kanał monotoniczny: przy wysokim BDP presja fiskalno-dłużna zaczyna dominować nad pozytywnym impulsem popytowym. \\
\addlinespace
Koszt kapitału &
Inflacja, reakcja NBP, dług publiczny i premia rynkowa wpływają na koszt finansowania inwestycji firm. &
inflacja, stopa NBP, obligacje, dług korporacyjny, inwestycje/PKB &
Stopa NBP rośnie umiarkowanie, ale rentowności obligacji i koszt długu korporacyjnego rosną wyraźniej; inwestycje prywatne spadają. \\
\addlinespace
Sektor zewnętrzny &
Wyższy popyt i importochłonność inwestycji wpływają na import, kurs PLN/EUR i rachunek bieżący. &
kurs, import, rachunek bieżący, ceny importu &
PLN słabnie, import rośnie, a rachunek bieżący pogarsza się; część impulsu popytowego odpływa przez kanał zewnętrzny. \\
\addlinespace
Kredyt i banki &
Finansowanie inwestycji zależy od jakości portfela kredytowego, płynności i adekwatności kapitałowej banków. &
NPL, CAR, LCR, NSFR, upadłości banków &
Nie występuje kryzys bankowy; ograniczenie automatyzacji przy wysokim BDP wynika raczej z kosztu i opłacalności finansowania niż z upadłości banków. \\
\addlinespace
Automatyzacja i firmy &
Firmy porównują koszt pracy, popyt, koszt kapitału i dostępność finansowania; automatyzacja wymaga sfinansowania CapEx. &
adopcja AI/hybrid, inwestycje, nowe firmy i firmy kończące działalność &
Adopcja AI/hybrid rośnie do \texttt{--bdp 2000}, po czym spada; więcej transferu nie oznacza większej automatyzacji, jeśli spada zdolność finansowania inwestycji. \\
\bottomrule
\end{tabularx}
\end{table}
